\subsection{舒尔正交关系}

我们保留上一小节的记号。设 \(\lambda\) 是 \(\widehat{G}\) 的一个元素，\(u\) 与 \(v\) 是 \(\mathrm{End}_{\mathrm{K}}(\mathrm{V}_{\lambda})\) 的元素；由公式 (10)，我们有

\[
\hat {\tau} (j _ {\lambda} (u) j _ {\lambda} (v)) = d _ {\lambda} \operatorname{Tr} (u v).
\]

由于 \(\tau = \hat{\tau} \circ \overline{\mathcal{F}}\)，我们从公式 (2) 与 (15) 导出关系式

\[
d _ {\lambda} ^ {2} | \mathrm{G} | ^ {- 1} \sum_ {g \in \mathrm{G}} \operatorname{Tr} (u \pi_ {\lambda} (g)) \operatorname{Tr} (v \pi_ {\lambda} (g ^ {- 1})) = d _ {\lambda} \operatorname{Tr} (u v).
\]

我们先前已指出，在域 K 中 \(d_{\lambda}.1 \neq 0\)；由此得出

\[
| \mathrm{G} | ^ {- 1} \sum_ {g \in \mathrm{G}} \operatorname{Tr} (u \pi_ {\lambda} (g)) \operatorname{Tr} (v \pi_ {\lambda} (g ^ {- 1})) = d _ {\lambda} ^ {- 1} \operatorname{Tr} (u v).\tag{20}
\]

把这一关系式应用于 u 与 v 的秩 \(\leqslant 1\) 的情形；我们得到

\[
| \mathrm{G} | ^ {- 1} \sum_ {g \in \mathrm{G}} \left\langle x ^ {*}, \pi_ {\lambda} (g) x \right\rangle \left\langle y ^ {*}, \pi_ {\lambda} \left(g ^ {- 1}\right) y \right\rangle = d _ {\lambda} ^ {- 1} \left\langle x ^ {*}, y \right\rangle \left\langle y ^ {*}, x \right\rangle\tag{21}
\]

其中 \(x, y\) 属于 \(\mathrm{V}_{\lambda}\)，\(x^{*}, y^{*}\) 属于对偶 \(\mathrm{V}_{\lambda}^{*}\)（\(\mathrm{V}_{\lambda}\) 的对偶）。

对每个 \(\lambda \in \widehat{\mathbf{G}}\)，设 \((e_{\lambda,j})_{1\leqslant j\leqslant d_{\lambda}}\) 是 \(\mathrm{V}_{\lambda}\) 的一组基；用 \((\pi_{ij}^{\lambda}(g))\) 表示自同态 \(\pi_{\lambda}(g)\)（\(\mathrm{V}_{\lambda}\) 的自同态）关于这组基的矩阵。若用 \((e_{\lambda,i}^{*})_{1\leqslant i\leqslant d_{\lambda}}\) 表示 \(\mathrm{V}_{\lambda}^{*}\) 的与 \((e_{\lambda,j})\) 对偶的基，则有 \(\pi_{ij}^{\lambda}(g) = \langle e_{\lambda,i}^{*}, \pi_{\lambda}(g)e_{\lambda,j}\rangle\)，因此

\[
| \mathrm{G} | ^ {- 1} \sum_ {g \in \mathrm{G}} \pi_ {i j} ^ {\lambda} (g) \pi_ {k \ell} ^ {\lambda} (g ^ {- 1}) = d _ {\lambda} ^ {- 1} \delta_ {i \ell} \delta_ {j k}.\tag{22}
\]

现设 \(\lambda\) 与 \(\mu\) 是 \(\widehat{\mathbf{G}}\) 的两个不同元素，且 \(u \in \mathrm{End}_{\mathrm{K}}(\mathrm{V}_{\lambda})\)，\(v \in \mathrm{End}_{\mathrm{K}}(\mathrm{V}_{\mu})\)。同样由关系式 (15)，我们有

\[
\sum_ {g \in \mathrm{G}} \operatorname{Tr} (u \pi_ {\lambda} (g)) \operatorname{Tr} (v \pi_ {\mu} (g ^ {- 1})) = 0.\tag{23}
\]

如上所述，我们导出

\[
\sum_ {g \in \mathrm{G}} \langle x ^ {*}, \pi_ {\lambda} (g) x \rangle \langle y ^ {*}, \pi_ {\mu} (g ^ {- 1}) y \rangle = 0\tag{24}
\]

其中 \(x\in \mathrm{V}_{\lambda},x^{*}\in \mathrm{V}_{\lambda}^{*},y\in \mathrm{V}_{\mu}\)，以及 \(y^{*}\in \mathrm{V}_{\mu}^{*}\)。我们还有

\[
\sum_ {g \in \mathrm{G}} \pi_ {i j} ^ {\lambda} (g) \pi_ {k \ell} ^ {\mu} (g ^ {- 1}) = 0\tag{25}
\]

对 \(i,j\)（在 \([1,d_{\lambda}]\) 中）与 \(k,\ell\)（在 \([1,d_{\mu}]\) 中）成立。

关系式 (20) 至 (25) 称为舒尔正交关系。

注. —— 我们把代数 \(\operatorname{End}_{\mathrm{K}}(\mathrm{V}_{\lambda})\) 与矩阵代数 \(\mathbf{M}_{d_{\lambda}}(\mathrm{K})\) 等同，借助基 \((e_{\lambda,j})\)（\(V_{\lambda}\) 的基）。映射 \(\overline{F}^{-1}\) 是从代数 \(\prod_{\lambda}\mathbf{M}_{d_{\lambda}}(\mathrm{K})\) 到代数 K{[}G{]} 的一个同构。对 \(\mu\in\widehat{G}\)，用 \(E_{ij}^{\mu}\) 表示 \(\prod_{\lambda}\mathbf{M}_{d_{\lambda}}(\mathrm{K})\) 中这样的元素：其指标 \(\mu\) 的分量是矩阵单位 \(E_{ij}\)（\(\mathbf{M}_{d_{\mu}}(\mathrm{K})\) 的矩阵单位）（II, §10, No.~3, 第 341 页），其余分量均为零；令 \(u_{ij}^{\mu}=\overline{\mathcal{F}}^{-1}(\mathrm{E}_{ij}^{\mu})\)。元素族 \(u_{ij}^{\lambda}\)（\(\lambda\in\widehat{G}\)，\(1\leqslant i\leqslant d_{\lambda}\)，\(1\leqslant j\leqslant d_{\lambda}\)）是代数 K{[}G{]} 的一组基；其乘法表为

\[
u _ {i j} ^ {\lambda} u _ {k \ell} ^ {\mu} = \delta_ {\lambda \mu} \delta_ {j k} u _ {i \ell} ^ {\lambda}.\tag{26}
\]

此外，由公式 (15)，我们有

\[
u _ {i j} ^ {\lambda} = | \mathrm{G} | ^ {- 1} d _ {\lambda} \sum_ {g \in \mathrm{G}} \pi_ {j i} ^ {\lambda} (g ^ {- 1}) g.\tag{27}
\]

